% \usepackage{booktabs,threeparttable} required

\begin{table}[t]
\centering
\begin{threeparttable}
\caption{Effect of requested pool size ($\alpha{\times}k$, the Dinkelbach OUTPUT size) on Stage 1
iteration count, at fixed retrieval breadth (\texttt{top\_n}, which sets the INPUT size
$|D_{union}|$), tier\_10k, $\sigma=0.60$.}
\label{tab:wdc-dinkelbach-iters-by-output-tier10k}
\begin{tabular}{rrrrr}
\toprule
\texttt{top\_n} & Iters at $\alpha{\times}k$=50\tnote{a} & Iters at $\alpha{\times}k$=1000\tnote{a} & Mean over 20 cells\tnote{b} & Ratio\tnote{c} \\
\midrule
 500 & 2.20 & 0.53 & 0.98 & 4.12$\times$ \\
1000 & 2.27 & 0.63 & 1.42 & 3.58$\times$ \\
2000 & 2.33 & 0.63 & 1.87 & 3.68$\times$ \\
3000 & 2.30 & 1.23 & 2.11 & 1.86$\times$ \\
4000 & 2.37 & 1.77 & 2.18 & 1.34$\times$ \\
5000 & 2.37 & 2.00 & 2.24 & 1.18$\times$ \\
\bottomrule
\end{tabular}
\begin{tablenotes}
\footnotesize
\item[a] Mean Dinkelbach iterations over the 30 queries at the smallest ($k$=10, $\alpha$=5) and
  largest ($k$=50, $\alpha$=20) $\alpha{\times}k$ cells of the 5$\times$4 $(k,\alpha)$ grid used
  throughout \texttt{topn\_alpha\_k\_sweep.py} -- the two extremes of \texttt{topn\_vs\_dunion\_tables.tex}'s
  "20 $\alpha{\times}k$ cells" note. Per-iteration cost does not depend on this column (each
  iteration is one $O(m\log m)$ sort over the fixed input $D$); iteration COUNT does, and that
  dependence is what this table isolates.
\item[b] Mean over all 20 $(k,\alpha)$ cells at this \texttt{top\_n} -- the same quantity
  \texttt{topn\_vs\_dunion\_tables.tex}'s "Mean iters" column reports.
\item[c] Ratio of column [a]'s two values (smallest-pool iterations / largest-pool iterations).
  \textbf{Shrinks from 4.1$\times$ at \texttt{top\_n}=500 to 1.2$\times$ at \texttt{top\_n}=5000} --
  the justification for reporting a single $(k,\alpha)$-averaged mean per \texttt{top\_n} in
  \texttt{topn\_vs\_dunion\_tables.tex}: at small \texttt{top\_n}, $|D_{union}|$ is small enough
  that even the grid's smallest tested pool (50) sits close to it, so Dinkelbach nears saturation
  (fast convergence) at large pools but not small ones, and the pool-size effect dominates -- the
  20-cell mean there blends two genuinely different convergence regimes. At \texttt{top\_n}=5000,
  $|D_{union}|$ (mean $\approx$1703, see \texttt{topn\_vs\_dunion\_tables.tex}) is large enough
  that NONE of the tested pools (50--1000) come close to saturating it, so iteration count is
  governed almost entirely by \texttt{top\_n} (input size) regardless of which $(k,\alpha)$ cell is
  picked -- the mean is a stable, representative summary there, not an average over disparate
  regimes. This is why \texttt{top\_n}=5000 was chosen as the fixed operating point for every later
  table (\texttt{wdc\_index\_cost\_table.tex}, \texttt{wdc\_ilp\_scaling\_tables.tex}).
\item[d] \textbf{Regenerated 2026-08-14} from a fresh sweep after fixing \texttt{HnswRetriever} to
  index categorical attributes only (\texttt{experiments/context.py::categorical\_only\_vectors});
  numbers here supersede an earlier version of this table computed before that fix.
\end{tablenotes}
\end{threeparttable}
\end{table}

\begin{table}[t]
\centering
\begin{threeparttable}
\caption{Effect of requested pool size ($\alpha{\times}k$, the Dinkelbach OUTPUT size) on Stage 1
iteration count, at fixed retrieval breadth (\texttt{top\_n}, which sets the INPUT size
$|D_{union}|$), tier\_100k, $\sigma=0.60$.}
\label{tab:wdc-dinkelbach-iters-by-output-tier100k}
\begin{tabular}{rrrrr}
\toprule
\texttt{top\_n} & Iters at $\alpha{\times}k$=50\tnote{a} & Iters at $\alpha{\times}k$=1000\tnote{a} & Mean over 20 cells\tnote{b} & Ratio\tnote{c} \\
\midrule
 500 & 2.03 & 0.38 & 0.87 & 5.36$\times$ \\
1000 & 2.07 & 0.37 & 1.27 & 5.64$\times$ \\
2000 & 2.13 & 0.37 & 1.77 & 5.82$\times$ \\
3000 & 2.17 & 0.93 & 1.99 & 2.32$\times$ \\
4000 & 2.17 & 1.70 & 2.10 & 1.27$\times$ \\
5000 & 2.17 & 1.90 & 2.15 & 1.14$\times$ \\
\bottomrule
\end{tabular}
\begin{tablenotes}
\footnotesize
\item[a] Mean Dinkelbach iterations over the 30 queries at the smallest ($k$=10, $\alpha$=5) and
  largest ($k$=50, $\alpha$=20) $\alpha{\times}k$ cells of the 5$\times$4 $(k,\alpha)$ grid used
  throughout \texttt{topn\_alpha\_k\_sweep.py} -- the two extremes of \texttt{topn\_vs\_dunion\_tables.tex}'s
  "20 $\alpha{\times}k$ cells" note. Per-iteration cost does not depend on this column (each
  iteration is one $O(m\log m)$ sort over the fixed input $D$); iteration COUNT does, and that
  dependence is what this table isolates.
\item[b] Mean over all 20 $(k,\alpha)$ cells at this \texttt{top\_n} -- the same quantity
  \texttt{topn\_vs\_dunion\_tables.tex}'s "Mean iters" column reports.
\item[c] Ratio of column [a]'s two values (smallest-pool iterations / largest-pool iterations).
  \textbf{Shrinks from 5.8$\times$ at \texttt{top\_n}=2000 to 1.1$\times$ at \texttt{top\_n}=5000}
  (peaks at \texttt{top\_n}=2000 rather than 500 here, since tier\_100k's $|D_{union}|$ grows more
  slowly with \texttt{top\_n} at the low end -- see \texttt{topn\_vs\_dunion\_tables.tex}'s note
  that tier\_100k's mean $|D_{union}|\approx$6378 at \texttt{top\_n}=5000 is $\sim$3.7x tier\_10k's)
  -- same justification as tier\_10k: at \texttt{top\_n}=5000, $|D_{union}|$ is large enough that
  none of the tested pools (50--1000) approach saturation, so the 20-cell mean is a stable summary
  of \texttt{top\_n}'s effect alone, not a blend of different convergence regimes.
\item[d] \textbf{Regenerated 2026-08-14} from a fresh sweep after fixing \texttt{HnswRetriever} to
  index categorical attributes only (\texttt{experiments/context.py::categorical\_only\_vectors});
  numbers here supersede an earlier version of this table computed before that fix.
\end{tablenotes}
\end{threeparttable}
\end{table}
